\documentclass{beamer}
\usetheme{Madrid}
\usecolortheme{dolphin}
\definecolor{charcoal}{RGB}{54,69,79}
\setbeamercolor{structure}{fg=charcoal}
\setbeamercolor{palette primary}{bg=charcoal,fg=white}
\setbeamercolor{palette secondary}{bg=charcoal!90,fg=white}
\setbeamercolor{palette tertiary}{bg=charcoal!80,fg=white}
\setbeamercolor{palette quaternary}{bg=charcoal!70,fg=white}
\setbeamercolor{frametitle}{bg=charcoal,fg=white}
\setbeamercolor{title}{bg=charcoal,fg=white}
\usepackage{hyperref}
\usepackage{graphicx}
\usepackage{booktabs}
\usepackage{listings}
\usepackage{color}
\definecolor{dkgreen}{rgb}{0,0.6,0}
\definecolor{gray}{rgb}{0.5,0.5,0.5}
\definecolor{mauve}{rgb}{0.58,0,0.82}
\lstset{
  language=Java,
  aboveskip=3mm,
  belowskip=3mm,
  showstringspaces=false,
  columns=flexible,
  basicstyle={\small\ttfamily},
  numbers=none,
  keywordstyle=\color{blue},
  commentstyle=\color{dkgreen},
  stringstyle=\color{mauve},
  breaklines=true,
  breakatwhitespace=true,
  tabsize=2
}
\title[Inheritance]{Inheritance and Polymorphism}
\subtitle{COP2250: Java Programming}
\author{Kevin Pyatt, Ph.D.}
\institute{State College of Florida \\ Pyatt Labs}
\date{Week 13}
\begin{document}
\begin{frame}
  \titlepage
\end{frame}
\begin{frame}{Today's Objectives}
\begin{itemize}
  \item Extend a superclass using \texttt{extends}
  \item Call the parent constructor using \texttt{super()}
  \item Override methods with \texttt{@Override}
  \item Understand method lookup order: subclass first, then superclass
  \item Apply Heron's formula to compute triangle area
  \item Build \texttt{Triangle extends GeometricObject}
\end{itemize}
\end{frame}
\begin{frame}{Is-A vs Has-A}
\textbf{Before you write \texttt{extends}, ask the question:}
\vspace{0.3cm}
\begin{tabular}{lll}
\toprule
\textbf{Relationship} & \textbf{Example} & \textbf{Use} \\
\midrule
Is-A & Triangle is a GeometricObject & \texttt{extends} \\
Has-A & GameEngine has a World & field (composition) \\
\bottomrule
\end{tabular}
\vspace{0.3cm}
\textit{Get this backwards and your hierarchy fights you for the rest of the project.}
\end{frame}
\begin{frame}[fragile]{Extending a Class}
\begin{lstlisting}
public class Triangle extends GeometricObject {
    private double side1;
    private double side2;
    private double side3;
}
\end{lstlisting}
\vspace{0.3cm}
\texttt{Triangle} automatically inherits:
\begin{itemize}
  \item \texttt{color} and \texttt{filled} fields (private --- use getters)
  \item \texttt{getColor()}, \texttt{setColor()}, \texttt{isFilled()}, \texttt{setFilled()}
  \item \texttt{toString()} --- can be overridden
\end{itemize}
\vspace{0.2cm}
\textit{You get everything. You add what's specific to Triangle.}
\end{frame}
\begin{frame}[fragile]{super() in Constructors}
\textbf{The parent must be constructed before the child.}
\begin{lstlisting}
public Triangle(double side1, double side2, double side3) {
    super();          // calls GeometricObject constructor
    this.side1 = side1;
    this.side2 = side2;
    this.side3 = side3;
}
\end{lstlisting}
\vspace{0.3cm}
\begin{itemize}
  \item \texttt{super()} must be the \textbf{first line} of the constructor
  \item If you omit it, Java inserts a no-arg \texttt{super()} automatically
  \item If the parent has no no-arg constructor --- it crashes
\end{itemize}
\vspace{0.2cm}
\textit{Always write it explicitly. Never rely on the implicit call.}
\end{frame}
\begin{frame}[fragile]{@Override}
\textbf{GeometricObject has toString(). Triangle wants its own.}
\begin{lstlisting}
@Override
public String toString() {
    return super.toString()
         + "\nSide 1: " + side1
         + "\nSide 2: " + side2
         + "\nSide 3: " + side3;
}
\end{lstlisting}
\vspace{0.3cm}
\begin{itemize}
  \item \texttt{@Override} tells the compiler you intend to override
  \item If you typo the method name, the compiler catches it
  \item \texttt{super.toString()} calls the parent version first
\end{itemize}
\vspace{0.2cm}
\textit{Always call super.toString(). Never throw away the parent's work.}
\end{frame}
\begin{frame}[fragile]{Method Lookup Order}
\textbf{Java looks in the actual object's class first.}
\begin{lstlisting}
GeometricObject g = new Triangle(3, 4, 5);
System.out.println(g.toString());
// calls Triangle's toString() -- not GeometricObject's
\end{lstlisting}
\vspace{0.3cm}
\begin{tabular}{ll}
Reference type & \texttt{GeometricObject} \\
Actual object & \texttt{Triangle} \\
Method called & \texttt{Triangle.toString()} \\
\end{tabular}
\vspace{0.3cm}
\textit{The reference type does not determine which method runs. The object does.}\\
\textit{This is polymorphism.}
\end{frame}
\begin{frame}[fragile]{Heron's Formula}
\textbf{Area of a triangle from three sides --- no height needed.}
\begin{lstlisting}
public double getArea() {
    double s = (side1 + side2 + side3) / 2;
    return Math.sqrt(s * (s - side1)
                       * (s - side2)
                       * (s - side3));
}
\end{lstlisting}
\vspace{0.3cm}
\textbf{Example: sides 3, 4, 5}
\begin{tabular}{ll}
s & = (3 + 4 + 5) / 2 = 6 \\
area & = $\sqrt{6 \times 3 \times 2 \times 1}$ = $\sqrt{36}$ = 6.0 \\
\end{tabular}
\vspace{0.3cm}
\textit{Verify your formula with a 3-4-5 triangle. Area must be exactly 6.0.}
\end{frame}
\begin{frame}{Key Terms}
\begin{description}
  \item[Inheritance] A subclass acquires fields and methods of its superclass
  \item[\texttt{extends}] Keyword that creates an inheritance relationship
  \item[\texttt{super()}] Calls the parent class constructor
  \item[\texttt{@Override}] Declares intent to replace a parent method
  \item[Polymorphism] A parent reference can point to a child object; child method runs
  \item[Method lookup] Java searches the actual object's class first
\end{description}
\end{frame}
\begin{frame}{Common Mistakes}
\begin{tabular}{ll}
\toprule
\textbf{Mistake} & \textbf{Fix} \\
\midrule
Accessing private parent field directly & Use the getter \\
Forgetting \texttt{super()} & Parent must be built first \\
\texttt{@Override} removed due to error & Fix the signature, keep \texttt{@Override} \\
\texttt{toString()} not calling \texttt{super.toString()} & Always call super \\
Wrong Heron's formula parentheses & Check each subtraction is wrapped \\
\bottomrule
\end{tabular}
\end{frame}
\begin{frame}{Lab --- Shape Hierarchy}
\textbf{Build a small inheritance hierarchy. Implement in order:}
\begin{enumerate}
  \item Read \texttt{Shape.java} --- do not modify it
  \item Create \texttt{Circle.java} extending \texttt{Shape}
  \item Create \texttt{Rectangle.java} extending \texttt{Shape}
  \item Override \texttt{toString()} in both --- call \texttt{super.toString()}
  \item Create \texttt{ShapeTest.java} --- instantiate and print both
  \item Compile all files, run, verify output
\end{enumerate}
\vspace{0.3cm}
\textbf{Demo Requirement:} Screen-share, compile live, explain what \texttt{super()} does.\\
\vspace{0.3cm}
\textbf{Next:} Assignment 11 --- Triangle extends GeometricObject (Liang 11.1)
\end{frame}
\end{document}
